% Generated from validated Article Draft Result. Do not hand-edit; revise the structured draft instead.
\section{巨大化を成立させる裏側――ZeROとGShard}
\label{pkg:systems-scale}
\noindent\textbf{frontier modelの規模は、parameter数だけでは決まらない。2020年のDeepSpeed / ZeROとGShardは、optimizer stateやmemoryの分割、distributed sharding、conditional computationといったsystems設計が、巨大modelを実際に訓練・実行するための独立した問題になったことを示す。scaleはmodel architectureとsystem architectureの共同問題へ広がった。}\autocite{src-56ccf2fa5ce7,src-f1047e7f654c}

DeepSpeed / ZeROが可視化したのは、large-model trainingでmodel parameter以外の状態もmemoryを占有し、単純にdeviceへ載せるだけでは規模拡大が止まるというsystems側の制約だった。ZeROはoptimizer state、gradient、parameterなどをdata-parallel process間で分割する方向を打ち出し、modelを大きくする能力がmemory placementとcommunication designに依存することを前景化した。\autocite{src-f1047e7f654c}


GShardは別の方向からscaleをsystems問題として扱った。論文は巨大なneural networkを複数deviceへ分割して実行するshardingの仕組みと、conditional computationを組み合わせる設計を示している。ここで重要なのは特定のparameter数の順位ではなく、model graphのどの部分をどのdeviceで計算し、どのparameterをtokenごとにactiveにするかがarchitecture choiceになった点である。\autocite{src-56ccf2fa5ce7}


conditional computationは、全parameterをすべてのinputで同じように使うdense scalingとは異なる発想を持つ。model capacityを増やしても、各tokenでactiveになる計算経路を限定できれば、parameter countとper-example computeを完全には同じ速度で増やさずに済む可能性がある。2020年の段階で、この発想がdistributed systemsと結び付いてscale設計の一つの系譜を形成していたことが重要である。\autocite{src-56ccf2fa5ce7}


ZeROのstate partitioningとGShardのsharding / conditional computationを並べると、2020年のscale競争がmodel-sizeの一本線ではなかったことが分かる。memory、communication、parallel execution、active computationという複数の資源制約を扱うsystems layerが、frontier-scale researchの成立条件になった。この視点は翌年のsparse MoE、memory offload、さらに大規模なdistributed trainingを理解する土台になる。\autocite{src-56ccf2fa5ce7,src-f1047e7f654c}


\begin{claimboundary}
ZeROやGShardが示したmemory効率、training scale、compute効率、model品質に関する評価はproject・author側の主張である。本稿はそれらのsystems設計が2020年の一次資料に存在し、scale成立条件として扱われたことを確認しているが、性能値や効率改善を独立再現したものではない。\autocite{src-56ccf2fa5ce7,src-f1047e7f654c}
\end{claimboundary}
